\documentclass[11pt]{article}
\usepackage[review]{acl}
\usepackage{times}
\usepackage{latexsym}
\usepackage[T1]{fontenc}
\usepackage[utf8]{inputenc}
\usepackage{amsmath}
\usepackage{microtype}

\title{A Neuroscience-Grounded Signature for Belief Imparting by Content:\\
Convergent but Modest Evidence across Text Persuasion, Real fMRI, and an In-Silico Encoder}

\author{Anonymous ACL submission}

\begin{document}
\maketitle

\begin{abstract}
Large language models can shift human attitudes, and human neuroscience shows that belief change and
value-based choice have predictable neural signatures. These literatures have not been joined. We ask
whether one neuroscience-grounded signature, computed from content, predicts real belief and choice,
generalizes across domains, holds on real fMRI, and can serve as an in-silico instrument that scores
content for belief-likelihood. We define the signature as $B = \text{engagement} + \text{value} -
\text{resistance}$ and test it in five stages on public data. A synthetic control validates the apparatus.
On ChangeMyView (3866 matched argument pairs), the engagement channel is higher for the argument that
changed a view (Wilcoxon $p = 9\mathrm{e}{-25}$), though the composite adds no pooled value over content
features. On Persuasion-for-Good, the effect points the same way and is significant when pooled
(coefficient 0.075, 95\% CI [0.036, 0.114]). On real fMRI, the signature computed from brain activity
predicts aggregate choice out-of-sample ($r = 0.182$, permutation $p = 0.003$), and a parsimonious
approach-minus-avoidance composite generalizes where richer models overfit. An in-silico encoder
generalizes across domains and isolates two robust levers (more evidence, less identity threat). The
effects are consistent and out-of-sample but modest. The contribution is a validated, interpretable
belief-scoring instrument, framed as a measurement and risk-characterization tool.
\end{abstract}

\section{Introduction}
Conversational AI can change human attitudes. Model-written messages shift political attitudes about as
much as human-written ones \citep{bai2025llm}, an effect driven mostly by post-training and prompting
rather than personalization or scale \citep{hackenburg2025levers, argyle2025testing}, and a meta-analysis
finds no mean difference between model and human persuasion but large heterogeneity, with study moderators
explaining most of the variance \citep{holbling2025meta}. This literature measures behavioral outcome and
carries no neural signal.

Human neuroscience supplies the missing half. Neural responses to persuasive messages predict later
behavior change beyond stated intention, with medial prefrontal cortex carrying the signal
\citep{falk2010predicting}. Nucleus accumbens activity forecasts aggregate choice, in several settings
better than choice itself \citep{genevsky2017when, knutson2018neuroforecasting, tong2020brain}. Belief
resistance appears in default-mode and dorsomedial prefrontal cortex \citep{kaplan2016neural}. Each is a
passive observation, and none asks whether such a signature, read from content, carries information that
content analysis misses.

We join the two. The research question is whether one neuroscience-grounded signature, computed from
content, predicts real belief and choice, generalizes across domains, holds on real neural data, and can
serve as an in-silico belief-scoring instrument. We answer it in five stages on public data. The aim is a
validated measurement instrument, not a deployed persuasion system.

\section{The Signature}
For a piece of content we read three quantities, each tied to a published result: engagement $E$ (medial
prefrontal cortex; \citealp{falk2010predicting}), value $V$ (nucleus accumbens;
\citealp{genevsky2017when}), and resistance $R$ (default-mode and anterior insula;
\citealp{kaplan2016neural}). The signature is $B = z(E) + z(V) - z(R)$. In the text arms $E, V, R$ are
computed from transparent lexical channels; in the neural arm they are read from real regional activity;
in the in-silico encoder they are intermediate variables constrained to the same structure.

\section{Data and Methods}
All data are public. ChangeMyView / Winning Arguments \citep{tan2016winning}, via ConvoKit: 19571 labeled
challenger arguments in 3866 matched pairs, outcome a delta awarded (view changed). Persuasion-for-Good
\citep{wang2019persuasion}: 1017 charity-donation dialogues, outcome the persuadee donated. A stock
neuroforecasting task \citep{kuhnen2005neural}: group-averaged BOLD in nucleus accumbens, medial
prefrontal cortex, and anterior insula, with aggregate buy fraction as outcome across two experiments.

Synthetic validation uses a stand-in with a known ground truth and three conditions (signal present,
content-only, scrambled). Text arms use within-pair sign tests, grouped cross-validation, bootstrap
intervals, per-channel decomposition, and a fixed-effect cross-corpus combination. The neural arm uses
cross-experiment out-of-sample prediction with permutation nulls and a six-model comparison. The in-silico
encoder is evaluated by cross-domain training, bootstrap factor variance, and counterfactual optimization.
All code and results will be released.

\section{Results}
\paragraph{Apparatus validity.} On the synthetic stand-in the instrument detects a planted signal
(delta-AUC $+0.067$), reports null when belief is content-only, and returns chance when the signature is
scrambled (AUC 0.497). Across twelve seeds the increment test detects the signal in 12 of 12
positive-control seeds and reports the correct null in 11 of 12 content-only seeds.

\paragraph{Real belief change.} On ChangeMyView, within matched pairs holding topic and audience fixed,
the engagement channel is higher for the argument that changed a view in 57.9\% of pairs (Wilcoxon $p =
9\mathrm{e}{-25}$) against a permutation null of 0.506. The value channel is null (0.487, $p = 0.07$); the
resistance text-proxy runs opposite to theory, so the composite is dominated by engagement. The signature
adds no pooled value over content features (content AUC 0.569, content plus $B$ 0.569), a null we read as
consistent with the meta-analytic heterogeneity.

\paragraph{Cross-domain generalization.} On Persuasion-for-Good, a different domain, the engagement
coefficient points the same way (ChangeMyView $+0.076$, significant; Persuasion-for-Good $+0.059$, same
sign, not individually significant), and the pooled estimate is significant ($+0.075$, 95\% CI [0.036,
0.114]). Value does not generalize.

\paragraph{Real neural data.} On real fMRI the signature predicts aggregate choice out-of-sample
(cross-experiment $r = 0.182$, AUC 0.571, permutation $p = 0.003$). The composite is the only predictor
sign-stable across experiments; single regions flip sign. The effect is near-null in one experiment
($r = 0.05$) and strong in the other ($r = 0.32$).

\paragraph{What generalizes best.} A six-model comparison, guided by the neuroforecasting and small-sample
regularization literatures \citep{genevsky2017when, tong2020brain, mortazavi2025deconstructing}, shows the
parsimonious approach-minus-avoidance composite (nucleus accumbens minus anterior insula) is the best
generalizer ($r = 0.180$, permutation $p = 0.002$). Every multivariate or learned model generalizes worse
(ridge on all features $r = 0.093$, not significant). The affect composite is positive across four
anticipatory time windows.

\paragraph{In-silico encoder.} A brain-structured encoder maps content factors through an
approach-minus-avoidance layer to a belief score \citep{wang2025foundation, walters2022predicting,
kriegeskorte2018interpreting}. Trained on one text corpus and applied to the other, the brain-state
read-out generalizes slightly better than the unconstrained model both ways (AUC 0.546 and 0.536 versus
0.537 and 0.519). Of six factors, only two are robust across both corpora: evidence raises belief and
identity threat lowers it. Counterfactually altering these two on low-belief content raises the predicted
belief score (ChangeMyView 0.600 to 0.625, Persuasion-for-Good 0.450 to 0.541).

\section{A Formal Model}
A companion model places the signature inside a dynamical account: belief is a state in a cusp potential
\citep{vandermaas2003sudden}, the persuasive force carries the signature term, gated by bounded confidence
\citep{hegselmann2002opinion} and asymmetric updating \citep{palminteri2022computational,
lefebvre2017behavioural}. Belief robustness is a computable functional (well stiffness and distance to the
fold, scaling as involvement to the three-halves). Four predictions are reproduced in the implementation:
hysteresis, a foot-in-the-door advantage of sequential in-band messaging, super-linear robustness scaling,
and susceptibility divergence at the fold. These are verified in code, not yet against human behavior.

\section{Discussion}
Three independent analyses converge on the same two levers. The ChangeMyView within-pair test, the real
fMRI composite, and the in-silico encoder all point to an approach-and-engagement signal minus an
avoidance-and-resistance signal. That the same structure carries a real, out-of-sample, permutation-
validated signal across text persuasion and real neural choice, and isolates the same interpretable levers
in an in-silico encoder, is the central finding. It is convergent but modest: an eight-point within-pair
lift, a cross-experiment $r$ near 0.18, a cross-domain AUC near 0.54. The signature is best used to score
content and characterize risk, not to forecast individual outcomes.

\section*{Limitations}
The text arms use lexical proxies for the neural quantity, not a learned brain encoder. The neural arm
uses real fMRI but a different behavior (financial choice) than belief change, and its effect is
experiment-dependent. The resistance channel runs opposite to theory in text. No single dataset closes the
loop from content to a neural state to a belief outcome in the same people, which does not yet exist
publicly at scale. The formal model's predictions are verified in code, not against human data. These bound
the claim to convergent, modest, out-of-sample evidence.

\section*{Ethics Statement}
This work develops and releases a scoring and risk-characterization instrument, not a deployed persuasion
system, and treats the optimization as a counterfactual sensitivity analysis rather than a generator
pointed at people. All data are public; no new human-subjects data are collected. The interpretable belief
levers support detection and inoculation research.

\bibliography{refs}

\end{document}
